\documentclass{amsart}
% For dealing with references we use the comment environment
\usepackage{verbatim}
\newenvironment{reference}{\comment}{\endcomment}
%\newenvironment{reference}{}{}
\newenvironment{slogan}{\comment}{\endcomment}
\newenvironment{history}{\comment}{\endcomment}

% For commutative diagrams we use Xy-pic
\usepackage[all]{xy}

% We use 2cell for 2-commutative diagrams.
\xyoption{2cell}
\UseAllTwocells

% We use multicol for the list of chapters between chapters
\usepackage{multicol}

% This is generally recommended for better output
\usepackage{lmodern}
\usepackage[T1]{fontenc}

% For cross-file-references

% Package for hypertext links:
\usepackage{hyperref}

% For any local file, say "hello.tex" you want to link to please

% Theorem environments.
%
\theoremstyle{plain}
\newtheorem{theorem}[subsection]{Theorem}
\newtheorem{proposition}[subsection]{Proposition}
\newtheorem{lemma}[subsection]{Lemma}

\theoremstyle{definition}
\newtheorem{definition}[subsection]{Definition}
\newtheorem{example}[subsection]{Example}
\newtheorem{exercise}[subsection]{Exercise}
\newtheorem{situation}[subsection]{Situation}

\theoremstyle{remark}
\newtheorem{remark}[subsection]{Remark}
\newtheorem{remarks}[subsection]{Remarks}

\numberwithin{equation}{subsection}

% Macros
%
\def\lim{\mathop{\mathrm{lim}}\nolimits}
\def\colim{\mathop{\mathrm{colim}}\nolimits}
\def\Spec{\mathop{\mathrm{Spec}}}
\def\Hom{\mathop{\mathrm{Hom}}\nolimits}
\def\Ext{\mathop{\mathrm{Ext}}\nolimits}
\def\SheafHom{\mathop{\mathcal{H}\!\mathit{om}}\nolimits}
\def\SheafExt{\mathop{\mathcal{E}\!\mathit{xt}}\nolimits}
\def\Sch{\mathit{Sch}}
\def\Mor{\mathop{\mathrm{Mor}}\nolimits}
\def\Ob{\mathop{\mathrm{Ob}}\nolimits}
\def\Sh{\mathop{\mathit{Sh}}\nolimits}
\def\NL{\mathop{N\!L}\nolimits}
\def\CH{\mathop{\mathrm{CH}}\nolimits}
\def\proetale{{pro\text{-}\acute{e}tale}}
\def\etale{{\acute{e}tale}}
\def\QCoh{\mathit{QCoh}}
\def\Ker{\mathop{\mathrm{Ker}}}
\def\Im{\mathop{\mathrm{Im}}}
\def\Coker{\mathop{\mathrm{Coker}}}
\def\Coim{\mathop{\mathrm{Coim}}}

% Boxtimes
%
\DeclareMathSymbol{\boxtimes}{\mathbin}{AMSa}{"02}

%
% Macros for moduli stacks/spaces
%
\def\QCohstack{\mathcal{QC}\!\mathit{oh}}
\def\Cohstack{\mathcal{C}\!\mathit{oh}}
\def\Spacesstack{\mathcal{S}\!\mathit{paces}}
\def\Quotfunctor{\mathrm{Quot}}
\def\Hilbfunctor{\mathrm{Hilb}}
\def\Curvesstack{\mathcal{C}\!\mathit{urves}}
\def\Polarizedstack{\mathcal{P}\!\mathit{olarized}}
\def\Complexesstack{\mathcal{C}\!\mathit{omplexes}}
% \Pic is the operator that assigns to X its picard group, usage \Pic(X)
% \Picardstack_{X/B} denotes the Picard stack of X over B
% \Picardfunctor_{X/B} denotes the Picard functor of X over B
\def\Pic{\mathop{\mathrm{Pic}}\nolimits}
\def\Picardstack{\mathcal{P}\!\mathit{ic}}
\def\Picardfunctor{\mathrm{Pic}}
\def\Deformationcategory{\mathcal{D}\!\mathit{ef}}

\setlength{\emergencystretch}{3em}
\begin{document}
\title[Stacks Project, Tag 0E8X]{575 --- Tricanonical very ampleness and first-cohomology vanishing for proper nodal curves of genus at least two with ample dualizing sheaf and ground-field global functions}
\maketitle
\noindent Copyright (C) 2005--2025 Johan de Jong. Stacks Project authors. Source: \href{https://stacks.math.columbia.edu/tag/0E8X}{Tag 0E8X}.

\noindent Editorial difficulty note (not author text). Difficulty H1: The proof must separate all degree-two subschemes even on reducible nodal curves. It reduces the obstruction to numerical inequalities on every connected subcurve, excludes rational tails and bridges using canonical ampleness, and handles the whole curve separately by Riemann\textendash{}Roch..

\noindent Editorial provenance note (not author text). History: excerpt from revision \texttt{a0444\allowbreak 6e57e\allowbreak c1fbc\allowbreak 252a8\allowbreak 71afc\allowbreak ec775\allowbreak 2fb28\allowbreak 07b14}, curves.tex, lines 6573--6635. Reference presentation was adapted; original-author context and core support blocks were retained or added. The mathematical wording is unchanged. Licensed under GNU FDL 1.2 or later; no invariant sections or cover texts. See ../licenses/COPYING. Context blocks reproduce author definitions and supporting results; they are not additional counted problems.

\begin{lemma}
\label{lemma-duality-dim-1}
Let $X$ be a proper scheme of dimension $\leq 1$ over a field $k$.
There exists a dualizing complex $\omega_X^\bullet$ with the
following properties
\begin{enumerate}
\item $H^i(\omega_X^\bullet)$ is nonzero only for $i = -1, 0$,
\item $\omega_X = H^{-1}(\omega_X^\bullet)$
is a coherent Cohen-Macaulay module whose support is the
irreducible components of dimension $1$,
\item for $x \in X$ closed, the module $H^0(\omega_{X, x}^\bullet)$
is nonzero if and only if either
\begin{enumerate}
\item $\dim(\mathcal{O}_{X, x}) = 0$ or
\item $\dim(\mathcal{O}_{X, x}) = 1$
and $\mathcal{O}_{X, x}$ is not Cohen-Macaulay,
\end{enumerate}
\item for $K \in D_\QCoh(\mathcal{O}_X)$ there are functorial
isomorphisms\footnote{This property
characterizes $\omega_X^\bullet$ in $D_\QCoh(\mathcal{O}_X)$
up to unique isomorphism by the Yoneda lemma. Since $\omega_X^\bullet$
is in $D^b_{\textit{Coh}}(\mathcal{O}_X)$ in fact it suffices to consider
$K \in D^b_{\textit{Coh}}(\mathcal{O}_X)$.}
$$
\Ext^i_X(K, \omega_X^\bullet) = \Hom_k(H^{-i}(X, K), k)
$$
compatible with shifts and distinguished triangles,
\item there are functorial isomorphisms
$\Hom(\mathcal{F}, \omega_X) = \Hom_k(H^1(X, \mathcal{F}), k)$
for $\mathcal{F}$ quasi-coherent on $X$,
\item if $X \to \Spec(k)$ is smooth of relative dimension $1$,
then $\omega_X \cong \Omega_{X/k}$.
\end{enumerate}
\end{lemma}

\begin{lemma}
\label{lemma-tricanonical}
Let $k$ be a field. Let $X$ be a proper scheme over $k$ of dimension $1$ with
$H^0(X, \mathcal{O}_X) = k$ having genus $g \geq 2$. Assume the singularities
of $X$ are at-worst-nodal and $\omega_X$ is ample. Then
$\omega_X^{\otimes 3}$ is very ample and $H^1(X, \omega_X^{\otimes 3}) = 0$.
\end{lemma}

\begin{proof}
Combining Varieties, Lemma
\href{https://stacks.math.columbia.edu/tag/0B5Y}{lemma ampleness in terms of degrees components (Tag 0B5Y)} and
Lemmas \href{https://stacks.math.columbia.edu/tag/0E63}{lemma rational tail negative (Tag 0E63)} and \href{https://stacks.math.columbia.edu/tag/0E64}{lemma rational bridge zero (Tag 0E64)}
we see that $X$ contains no rational tails or bridges.
Then we see that $\omega_X^{\otimes 3}$ is globally generated
by Lemma \href{https://stacks.math.columbia.edu/tag/0E3L}{lemma contracting rational tails (Tag 0E3L)}.
Choose a $k$-basis $s_0, \ldots, s_n$ of
$H^0(X, \omega_X^{\otimes 3})$. We get a morphism
$$
\varphi_{\omega_X^{\otimes 3}, (s_0, \ldots, s_n)} :
X \longrightarrow \mathbf{P}^n_k
$$
See Constructions, Section \href{https://stacks.math.columbia.edu/tag/01ND}{section projective space (Tag 01ND)}.
The lemma asserts that this morphism is a closed immersion.
To check this we may replace $k$ by its algebraic closure, see
Descent, Lemma \href{https://stacks.math.columbia.edu/tag/02L6}{lemma descending property closed immersion (Tag 02L6)}.
Thus we may assume $k$ is algebraically closed.

\medskip\noindent
Assume $k$ is algebraically closed.
We will use Varieties, Lemma
\href{https://stacks.math.columbia.edu/tag/0E8T}{lemma variant separate points tangent vectors (Tag 0E8T)}
to prove the lemma.
Let $Z \subset X$ be a closed subscheme of degree $2$ over $Z$
with ideal sheaf $\mathcal{I} \subset \mathcal{O}_X$.
We have to show that
$$
H^0(X, \mathcal{L}) \to H^0(Z, \mathcal{L}|_Z)
$$
is surjective. Thus it suffices to show that
$H^1(X, \mathcal{I}\mathcal{L}) = 0$.
To do this we will use Lemma \href{https://stacks.math.columbia.edu/tag/0E3E}{lemma vanishing on gorenstein (Tag 0E3E)}.
Thus it suffices to show that
$$
3\deg(\omega_X|_Y) > -2\chi(Y, \mathcal{O}_Y) + \deg(Z \cap Y)
$$
for every reduced connected closed subscheme $Y \subset X$.
Since $k$ is algebraically closed and $Y$ connected and reduced
we have $H^0(Y, \mathcal{O}_Y) = k$ (Varieties, Lemma
\href{https://stacks.math.columbia.edu/tag/0BUG}{lemma proper geometrically reduced global sections (Tag 0BUG)}).
Hence $\chi(Y, \mathcal{O}_Y) = 1 - \dim H^1(Y, \mathcal{O}_Y)$.
Thus we have to show
$$
3\deg(\omega_X|_Y) > -2 + 2\dim H^1(Y, \mathcal{O}_Y) + \deg(Z \cap Y)
$$
which is true by Lemma \href{https://stacks.math.columbia.edu/tag/0E3J}{lemma no rational tail (Tag 0E3J)}
except possibly if $Y = X$ or if $\deg(\omega_X|_Y) = 0$.
Since $\omega_X$ is ample the second possibility does not
occur (see first lemma cited in this proof). Finally, if
$Y = X$ we can use Riemann-Roch (Lemma \href{https://stacks.math.columbia.edu/tag/0BS6}{lemma rr (Tag 0BS6)})
and the fact that $g \geq 2$ to see that the inquality holds.
The same argument with $Z = \emptyset$ shows that
$H^1(X, \omega_X^{\otimes 3}) = 0$.
\end{proof}
\end{document}
